%! Addition and Multiplication Rules — Unit 3, Lesson 3.3 — Slides (Beamer)
% GRADUAL RELEASE deck (5/45/5/5): title, targets, warm-up, hook (dark, UNRESOLVED), I DO divider,
% notes rows (crux flagged redacc), Guided Practice (LIVE, un-answered), debrief, You do
% (assigned, not started), close. Standards: PS.P.1a, c; AFDA.DA.3e, f.
\documentclass[aspectratio=169, 11pt]{beamer}
\usepackage{saar-beamer}   % provides \ceruleanheader{} and \sectionlabel[color]{}

\newcommand{\redink}[1]{{\color{redacc}\bfseries #1}}
\newcommand{\venn}[6]{%
  \begin{tikzpicture}[x=1cm, y=1cm]
    \draw[thick, charcoal] (0,0) rectangle (7,3);
    \fill[cerulean, fill opacity=0.18] (#1,1.45) circle[radius=1.15];
    \fill[goldacc, fill opacity=0.25] (#2,1.45) circle[radius=1.15];
    \draw[thick, cerulean] (#1,1.45) circle[radius=1.15];
    \draw[thick, goldacc] (#2,1.45) circle[radius=1.15];
    #3
    \node[font=\bfseries] at (6.4,0.4) {#4};
    \node[font=\small\bfseries] at (0.5,2.7) {#5};
    \node[font=\small\bfseries] at (6.5,2.7) {#6};
  \end{tikzpicture}}

\begin{document}

{
\setbeamercolor{background canvas}{bg=cerulean}
\begin{frame}[plain]
  \vspace{2.8cm}\hspace{0.55cm}%
  \begin{minipage}{0.9\textwidth}
    {\color{goldacc}\bfseries\small LESSON 3.3}\\[0.5cm]
    {\color{white}\bfseries\LARGE Addition and Multiplication Rules}\\[1.2cm]
    {\color{frostmid}\small Statistical Analysis \& Algebraic Reasoning~$\cdot$~Unit 3~$\cdot$~PS.P.1a, c~$\cdot$~AFDA.DA.3e, f}
  \end{minipage}
\end{frame}
}

\begin{frame}[t, plain]
  \ceruleanheader{Today's targets}
  \sectionlabel{Lesson 3.3}
  By the end of today, I can\ldots
  \begin{itemize}\setlength{\itemsep}{4pt}
    \item use the \textbf{addition rule} on a Venn diagram, subtracting the overlap.
    \item find ``at least one'' with the complement rule.
    \item use the \textbf{multiplication rule} on a tree for \textbf{independent} and
          \textbf{dependent events}.
    \item classify two events by computing --- and explain why mutually exclusive is not
          independent.
  \end{itemize}
  \begin{block}{How today works}
    Warm-up $\to$ \textbf{notes together} $\to$ \textbf{one context we work as a class} $\to$
    debrief $\to$ \textbf{the homework, assigned today and done outside class}.
  \end{block}
\end{frame}

\begin{frame}[t, plain]
  \ceruleanheader{Warm-Up --- 5 minutes, silent}
  \sectionlabel{back to Lesson 3.2 $\cdot$ the concert raffle}
  \vspace{0.3cm}
  \begin{enumerate}\setlength{\itemsep}{7pt}
    \item Compute: $0.6 + 0.5 - 0.3$, \quad $0.6 \cdot 0.5$, \quad $1 - 0.2$.
    \item Multiply: $\frac{3}{10} \cdot \frac{3}{10}$ \quad and \quad
          $\frac{3}{10} \cdot \frac{2}{9}$.
    \item Read the raffle Venn on your page. Why is the $S$ count plus the $E$ count more than 40?
  \end{enumerate}
  \vspace{0.4cm}
  \begin{center}
    {\color{cerulean}\bfseries Keep the page --- the notes open with this Venn.}
  \end{center}
\end{frame}

{
\setbeamercolor{background canvas}{bg=cerulean}
\begin{frame}[plain]
  \vspace{1.2cm}
  \begin{center}
    {\color{frostmid}\bfseries\small ONE WINNING TICKET. SOLD TO A NINTH GRADER, OR SOLD TO A TWELFTH GRADER. NEVER BOTH.}\\[0.8cm]
    {\color{white}\bfseries\Large ``They can't happen together, so they have nothing to do with each other. They must be independent.''}\\[0.8cm]
    {\color{goldacc}\bfseries\large Independent, dependent, or can't tell?}\\[0.5cm]
    {\color{frostmid}\small Vote, and write why. We vote again at problem 10.}
  \end{center}
\end{frame}
}

{
\setbeamercolor{background canvas}{bg=cerulean}
\begin{frame}[plain]
  \vspace{1.8cm}\hspace{0.8cm}%
  \begin{minipage}{0.86\textwidth}
    {\color{goldacc}\bfseries\small GUIDED NOTES \ \textbullet\ \ I DO}\\[0.7cm]
    {\color{white}\bfseries\Large Open to your Guided Notes page. We work the table row by row
    --- fill each term as we name it, not before.}\\[0.9cm]
    {\color{frostmid}\small Four terms go in the vocabulary box today.}
  \end{minipage}
\end{frame}
}

\begin{frame}[t, plain]
  \ceruleanheader{Subtract the overlap}
  \sectionlabel{notes row 1 $\cdot$ Addition Rule $\cdot$ problems 1--2}
  \begin{center}
    $P(A \text{ or } B) = P(A) + P(B) - P(A \text{ and } B)$\\[6pt]
    \venn{2.75}{4.25}{\node[font=\Large\bfseries] at (2.1,1.45) {12};
      \node[font=\Large\bfseries] at (3.5,1.45) {12};
      \node[font=\Large\bfseries] at (4.9,1.45) {8};}{8}{$S$}{$E$}
  \end{center}
  \begin{itemize}\setlength{\itemsep}{3pt}
    \item Problem 1: $.6 + .5 - .3 = \mathbf{.8}$; counting, $32$ of $40$.
    \item Problem 2: Jamal's $1.1$ \redink{cannot be a probability}: the 12 overlap tickets
          were counted twice.
  \end{itemize}
\end{frame}

\begin{frame}[t, plain]
  \ceruleanheader{Overlap of zero}
  \sectionlabel{notes row 2 $\cdot$ Mutually Exclusive $\cdot$ problems 3--4}
  \begin{center}
    \venn{1.9}{5.1}{\node[font=\Large\bfseries] at (1.9,1.45) {8};
      \node[font=\Large\bfseries] at (5.1,1.45) {10};}{22}{$N$}{$T$}
  \end{center}
  \begin{itemize}\setlength{\itemsep}{3pt}
    \item Mutually exclusive: $P(A \text{ and } B) = 0$, so $P(A \text{ or } B) = P(A) + P(B)$.
    \item Problem 3: $.2 + .25 = \mathbf{.45}$.
    \item Problem 4: $P(\text{at least one}) = 1 - P(\text{none}) = 1 - .2 = \mathbf{.8}$.
  \end{itemize}
\end{frame}

\begin{frame}[t, plain]
  \ceruleanheader{Ticket put back: the branches stay the same}
  \sectionlabel{notes row 3 $\cdot$ Trees $\cdot$ problems 5--6}
  \vspace{4pt}
  \textbf{Bowl:} 10 tickets, 3 gold and 7 white. Two draws.
  \begin{itemize}\setlength{\itemsep}{6pt}
    \item \textbf{Independent:} one event does not change the other's probability.
    \item Problem 5: $\frac{3}{10} \cdot \frac{3}{10} = \mathbf{\frac{9}{100}}$.
    \item Problem 6: Ana's $\frac{3}{10} + \frac{3}{10}$ counts (gold, gold) twice.
          \redink{$1 - \frac{7}{10} \cdot \frac{7}{10} = \frac{51}{100}$}.
  \end{itemize}
\end{frame}

\begin{frame}[t, plain]
  \ceruleanheader{Ticket kept out: the branches change}
  \sectionlabel{notes row 3 $\cdot$ Trees $\cdot$ problems 7--8}
  \vspace{4pt}
  \begin{itemize}\setlength{\itemsep}{6pt}
    \item After a gold is kept out: $\frac{2}{9}$ gold, $\frac{7}{9}$ white. After a white:
          $\frac{3}{9}$ gold, $\frac{6}{9}$ white.
    \item Problem 7: $\frac{3}{10} \cdot \frac{2}{9} = \frac{6}{90} = \mathbf{\frac{1}{15}}$.
    \item Problem 8: second gold after first gold is $\frac{3}{10}$ put back and $\frac{2}{9}$
          kept out.
  \end{itemize}
  \begin{block}{Multiplication rule}
    Independent: $P(A) \cdot P(B)$. \quad Dependent: the second factor is $B$'s probability
    \emph{after} $A$ happened.
  \end{block}
\end{frame}

\begin{frame}[t, plain]
  \ceruleanheader{Compute, then name}
  \sectionlabel{notes row 4 $\cdot$ Four Words $\cdot$ problem 9}
  \begin{columns}[t]
    \begin{column}{0.55\textwidth}
      \small
      \renewcommand{\arraystretch}{1.35}
      \begin{tabular}{@{} l l @{}}
      \textbf{Mutually exclusive} & $P(A \text{ and } B) = 0$ \\
      \textbf{Complementary} & also $P(A) + P(B) = 1$ \\
      \textbf{Independent} & $P(A \text{ and } B) = P(A)P(B)$ \\
      \textbf{Dependent} & $P(A \text{ and } B) \neq P(A)P(B)$ \\
      \end{tabular}
    \end{column}
    \begin{column}{0.41\textwidth}
      \begin{block}{Problem 9: $S$ and $E$}
        $P(S \text{ and } E) = .3$\\
        $P(S)P(E) = (.6)(.5) = .3$\\
        Equal: \redink{independent}, with an overlap.
      \end{block}
    \end{column}
  \end{columns}
\end{frame}

\begin{frame}[t, plain]
  \ceruleanheader{Can't happen together. So: dependent.}
  \sectionlabel[redacc]{notes row 4 $\cdot$ problems 10--11 $\cdot$ the whole point of today}
  \vspace{4pt}
  \begin{itemize}\setlength{\itemsep}{5pt}
    \item \textbf{Problem 10 --- vote again, then compute.}
    \item $P(N \text{ and } T) = \mathbf{0}$, \quad $P(N)P(T) = (.2)(.25) = \mathbf{.05}$.
          \redink{Not equal: dependent.}
    \item Once you know $N$ happened, the chance of $T$ drops from $.25$ to $0$.
    \item Problem 11: $S$ and $A$: $.6 + .4 = 1$, overlap $0$, product $.24$. Complementary and
          dependent.
  \end{itemize}
  \begin{block}{Mutually exclusive is not independent}
    When two events that each have some chance cannot happen together, knowing one happened
    tells you the other did not.
  \end{block}
\end{frame}

\begin{frame}[t, plain]
  \ceruleanheader{Guided Practice: Rain and the Bus}
  \sectionlabel[goldacc]{We do $\cdot$ problems 12--15 $\cdot$ you hold the pen}
  \begin{columns}[c]
    \begin{column}{0.42\textwidth}
      \scalebox{0.8}{\venn{2.75}{4.25}{\node[font=\bfseries] at (2.1,1.45) {.25};
        \node[font=\bfseries] at (3.5,1.45) {.15};
        \node[font=\bfseries] at (4.9,1.45) {.05};}{.55}{$R$}{$L$}}
    \end{column}
    \begin{column}{0.54\textwidth}
      \begin{itemize}\setlength{\itemsep}{4pt}
        \item[12.] Find $P(R \text{ or } L)$. What does it mean for Leah?
        \item[13.] Compute $P(R \text{ and } L)$ and $P(R)P(L)$. Independent or dependent?
        \item[14.] Two independent mornings: late both? Late at least once?
        \item[15.] $P(\text{early}) = .1$. Theo says late and early are independent. Answer him.
      \end{itemize}
    \end{column}
  \end{columns}
\end{frame}

\begin{frame}[t, plain]
  \ceruleanheader{Debrief: say it with two numbers}
  \sectionlabel{5 minutes $\cdot$ pens down, papers up --- we say these aloud}
  \vspace{4pt}
  \begin{enumerate}\setlength{\itemsep}{5pt}
    \item Problem 15, back on the board. Say the sentence: \emph{they can't happen together, so
          knowing one tells you the other did not.}
  \end{enumerate}
  \vspace{6pt}
  \begin{block}{Cold check --- hands down, thirty seconds}
    Flip one coin once. Event $H$: heads. Event $T$: tails. \textbf{Independent or dependent?
    Give me the two numbers.}
  \end{block}
\end{frame}

\begin{frame}[t, plain]
  \ceruleanheader{You do: the homework}
  \sectionlabel[goldacc]{Assigned today $\cdot$ on your own, outside class $\cdot$ due after two study halls}
  \begin{itemize}\setlength{\itemsep}{5pt}
    \item \textbf{On DeltaMath: \emph{Lesson 3.3: Addition and Multiplication Rules}.}
    \item The addition rule and the complement rule, in symbols.
    \item ``Or'' and ``and'' probabilities, and compound events on two stages.
    \item Independent or dependent: multiply, then compare --- from a table, and with and
          without replacement.
  \end{itemize}
  \begin{block}{DeltaMath: Lesson 3.3: Addition and Multiplication Rules. Scored 1 to 5, due at the start of the first class after two study halls.}
    Row 4 of your notes and the \emph{Keep in Mind} box are your reference.
  \end{block}
\end{frame}

{
\setbeamercolor{background canvas}{bg=cerulean}
\begin{frame}[plain]
  \vspace{1.5cm}\hspace{0.8cm}%
  \begin{minipage}{0.86\textwidth}
    {\color{goldacc}\bfseries\small BEFORE YOU GO}\\[0.7cm]
    {\color{white}\bfseries\Large Add and subtract the overlap. Multiply along the branches.
    And compute before you name: mutually exclusive is not independent.}\\[0.9cm]
    {\color{frostmid}\small Next: Lesson 3.4, Conditional Probability, Independence, and
    Simulation.}
  \end{minipage}
\end{frame}
}

\end{document}
